% ============================================================================
\section{Percepción: del archivo a los renglones}
\label{sec:percepcion}
% ============================================================================

La percepción es la frontera entre el mundo visual del documento y el mundo
simbólico del resto del sistema. Tiene dos responsabilidades, y conviene
separarlas desde el principio:

\begin{enumerate}
  \item \textbf{Leer}: extraer de cada página los renglones de la lista, con su
        descripción y su cantidad, respetando la estructura de tabla.
  \item \textbf{Decidir si el archivo es aceptable}: si corresponde a la lista
        de \emph{un solo} grado y si se leyó \emph{completo}.
\end{enumerate}

La primera se delega a un modelo multimodal. La segunda no se delega: es una
regla de negocio implementada de forma determinista sobre lo que el modelo
devuelve. Esta separación es la decisión de diseño más importante de la
etapa, y el resto de la sección la desarrolla.

\subsection{Definición del problema}

\begin{definicion}[Problema de percepción]
\label{def:percepcion}
\textbf{Entrada:} una cadena de bytes $A$ (un PDF o una imagen).
\textbf{Salida:} una tupla
\[
  \Phi_{\mathrm{per}}(A) = \bigl(\mathit{ok},\ G,\ P,\ I,\ \varepsilon\bigr)
\]
donde $\mathit{ok}\in\{0,1\}$ indica si el archivo se acepta, $G$ es la
secuencia de grados distintos detectados, $P$ el número de páginas,
$I=(\iota_1,\dots,\iota_m)$ la secuencia de ítems extraídos, cada uno
$\iota=(x,q^\circ,\gamma)$ con descripción $x$, cantidad opcional $q^\circ$ y
confianza opcional $\gamma$, y $\varepsilon$ un código de error que es nulo si
y solo si $\mathit{ok}=1$.

\noindent\textbf{Requisitos:}
\begin{enumerate}[label=(R\arabic*)]
  \item\label{req:completo} \emph{Completitud o rechazo.} Si
        $\mathit{ok}=1$, entonces $I$ contiene los ítems de \emph{todas} las
        páginas y ninguna página se reportó incompleta.
  \item\label{req:grado} \emph{Unicidad de grado.} Si $\mathit{ok}=1$,
        entonces $\card{G}\leq 1$.
  \item\label{req:nofab} \emph{No fabricación.} Ningún ítem de $I$ se
        produce fuera de la lectura del modelo, y un fallo del modelo nunca se
        traduce en ítems inventados.
\end{enumerate}
\end{definicion}

El requisito \ref{req:completo} existe por la asimetría de costos ya
mencionada. Sea $c_{\mathrm{FA}}$ el costo de \emph{aceptar} un archivo que no
debía aceptarse (una lista de otro grado o una lista truncada) y
$c_{\mathrm{FR}}$ el costo de \emph{rechazar} uno válido. Un falso rechazo le
cuesta a la familia volver a subir el archivo, con un mensaje que le dice qué
hacer. Una falsa aceptación produce un kit con los útiles equivocados que la
familia no puede detectar hasta que recibe el pedido. Como
$c_{\mathrm{FA}}\gg c_{\mathrm{FR}}$, toda regla ambigua de esta etapa se
inclina hacia el rechazo.

\subsection{Por qué leer páginas como imágenes}

La alternativa ingenua es extraer el texto del PDF y partirlo en líneas. Esa
alternativa falla de forma sistemática con las listas escolares, que casi
siempre son \emph{tablas}: una fila cuya descripción ocupa dos líneas
visuales se convierte en dos renglones, y la columna de cantidad queda
separada de su descripción. El error no es aleatorio sino estructural: aplanar
un objeto bidimensional a una secuencia unidimensional destruye la relación
fila-columna.

Por eso cada página se \emph{rasteriza} y se entrega como imagen a un modelo
multimodal, al que se le pide leerla \emph{como tabla}. La rasterización se
hace a 150~DPI, es decir, con un factor de escala $s=150/72\approx 2{,}08$
sobre las unidades de PDF (1/72 de pulgada). Una página carta queda en
aproximadamente $1275\times 1650$ píxeles. Esa resolución basta para que una
tabla de útiles se lea como tabla; aumentarla encarece cada llamada sin
cambiar lo que el modelo alcanza a ver. Una imagen suelta (JPG, PNG) ya es una
página y no se transforma.

\subsection{Lectura estructurada de una página}
\label{sec:per-lectura}

La lectura de la página $j$ es una función
\[
  \mathrm{Leer}(I_j, j) = (G_j,\ \kappa_j,\ \iota_{j,1},\dots,\iota_{j,m_j})
  \quad\text{o}\quad \textsc{Fallo},
\]
donde $G_j$ es la lista de grados que aparecen en la página y
$\kappa_j\in\{0,1\}$ indica si la página se leyó completa. Su implementación
tiene tres componentes:

\paragraph{Decodificación restringida por esquema.} La salida del modelo se
restringe a un esquema JSON con los campos \cod{grados} (arreglo de cadenas),
\cod{completa} (booleano) e \cod{items} (arreglo de objetos con
\cod{descripcion} obligatoria, \cod{cantidad} entera o nula y \cod{confianza}
numérica o nula). El esquema se pide en dos lugares a la vez: en el texto de la
instrucción, que es lo que el modelo usa para decidir cómo segmentar, y en la
configuración de la llamada, que es lo que el proveedor valida. La respuesta se
vuelve a validar localmente contra el mismo tipo; cualquier discrepancia
(campo faltante, tipo incorrecto, JSON mal formado) se convierte en
\textsc{Fallo}. La temperatura se fija en $0$, lo que equivale a decodificación
voraz: sobre la misma imagen, el modelo devuelve la secuencia más probable, lo
que reduce la variabilidad entre lecturas de la misma página.

\paragraph{Especificación de segmentación.} La instrucción al modelo define una
\emph{función de segmentación} de filas en ítems. Sus reglas, cada una
originada en un caso real del corpus, son:

\begin{enumerate}[label=(S\arabic*)]
  \item \emph{Una fila es un ítem}, aunque ocupe dos líneas visuales. No se
        corta una descripción a la mitad.
  \item \emph{La cantidad sale de la columna de cantidad}. Si no hay columna,
        el campo queda nulo; el modelo no la infiere.
  \item \emph{Los contenidos se separan.} «Cartuchera con: esferos, borrador»
        son tres ítems, contando la cartuchera.
  \item \emph{Las enumeraciones de colores se expanden.} «Pinturas acrílicas:
        cian, magenta y rojo» son tres ítems, uno por color, aunque los colores
        aparezcan al final de la frase.
  \item \emph{La cantidad total se reparte, no se repite.} Si la fila pide 4
        carpetas en 4 colores, cada ítem resultante lleva cantidad 1. Formalmente,
        si una fila con cantidad $Q$ se expande en $n$ ítems y $Q=n$, cada ítem
        recibe $1$, de modo que $\sum_\ell q_\ell = Q$ se conserva.
  \item \emph{No se separan las alternativas ni los productos enteros.}
        «Tajalápiz o portaminas» es un solo ítem; «Kit de 12 pinceles (planos,
        redondos)» también.
  \item \emph{Una serie de durezas es un solo ítem.} «Lápices HB, 2H, 2B, 6B»
        pide una serie, no cuatro lápices sueltos.
  \item \emph{Un uniforme es un ítem con su nombre}, sin enumerar sus prendas.
  \item \emph{Los grados se nombran}: «Cuarto», no «4°».
  \item \emph{Completitud.} Si la imagen está recortada, borrosa o ilegible,
        \cod{completa} es falso. Un texto que \emph{continúa} en la página
        siguiente no es una página cortada.
  \item \emph{No fabricación.} No se agregan ítems que no estén escritos.
\end{enumerate}

Las reglas S3 a S7 dan una tensión que la instrucción resuelve de forma
explícita. Separar demasiado convierte un producto en varios; separar poco
deja dos compras en un renglón. El criterio que las ordena es \emph{qué se
compra por aparte}: lo que se compra por aparte se separa, lo que se vende como
unidad se conserva.

\paragraph{Política de reintento.} Cada llamada tiene un tiempo límite de
$t_{\max}=40$~s. Los fallos se clasifican en \emph{pasajeros} (sin respuesta a
tiempo, error de red, códigos HTTP 429, 500, 502, 503, 504) y
\emph{permanentes} (cualquier otro código, como un 403 por credencial
inválida, o una respuesta que no cumple el esquema). Ante un fallo pasajero se
reintenta \emph{una} vez, después de una espera de 3~s. Ante uno permanente no
se reintenta, porque repetir la petición no cambia el resultado y solo
escondería el defecto. El peor caso por página es por tanto
$t_{\max}+3+t_{\max}=83$~s.

El beneficio del reintento crece con el número de páginas. Si cada llamada se
cuelga de forma independiente con probabilidad $p$, la probabilidad de que un
documento de $P$ páginas falle es
\[
  \Pr[\text{falla}]_{\text{sin reintento}} = 1-(1-p)^P,
  \qquad
  \Pr[\text{falla}]_{\text{con un reintento}} = 1-(1-p^2)^P.
\]
Como ilustración, con $p=0{,}05$ y $P=6$ la probabilidad baja de 26,5\,\%
a 1,5\,\%. El supuesto de independencia es optimista ---los fallos del
proveedor tienden a agruparse en el tiempo---, pero el efecto observado fue del
mismo signo: una página que se quedó colgada 60~s respondió en 6~s al pedirla
de nuevo. El efecto es importante porque, por el requisito
\ref{req:completo}, un solo fallo invalida el documento completo, y volver a
subirlo obliga a pagar la lectura de todas las páginas otra vez.

\subsection{Normalización de grados}
\label{sec:per-grados}

Para verificar el requisito \ref{req:grado} hay que contar cuántos grados
\emph{distintos} menciona el archivo. El modelo y el colegio escriben el mismo
grado de formas distintas («4», «4°», «Cuarto», «grado cuarto», «4 de
primaria»), así que primero se canonizan.

\begin{definicion}[Canonización de grado]
\label{def:canon-grado}
Sea $\mathcal{G}$ el conjunto de nombres canónicos de grado (Prejardín,
Jardín, Transición, Primero, \dots, Once) y $T\colon\Sigma^*\rightharpoonup
\mathcal{G}$ una tabla finita de equivalencias («4»$\mapsto$Cuarto,
«kinder»$\mapsto$Transición, «undecimo»$\mapsto$Once, \dots). La canonización
$\eta\colon\Sigma^*\to\Sigma^*$ se define en cuatro pasos:
\begin{enumerate}
  \item $y\gets$ minúsculas de $x$ sin diacríticos (descomposición NFD y
        eliminación de marcas combinantes), con los símbolos de ordinal
        $^{\circ}$ y $^{\mathrm{o}}$ reemplazados por espacios, todo carácter
        no alfanumérico reemplazado por espacio, los espacios colapsados y la
        palabra «grado» eliminada.
  \item Si $y\in\mathrm{dom}(T)$, $\eta(x)=T(y)$.
  \item Si la primera palabra $y_1$ de $y$ está en $\mathrm{dom}(T)$,
        $\eta(x)=T(y_1)$. Así «4 de primaria» y «cuarto de primaria» se
        identifican.
  \item En otro caso, $\eta(x)$ es $y$ con mayúscula inicial en cada palabra.
\end{enumerate}
El conjunto de grados distintos de un archivo es
$G=\mathrm{únicos}\bigl(\eta(g) : g\in G_1\cdots G_P,\ \eta(g)\neq\varepsilon\bigr)$,
en orden de primera aparición.
\end{definicion}

El paso 4 es deliberadamente \emph{conservador}. Un nombre de grado que la
tabla no reconoce no se descarta: se conserva limpio y cuenta como un grado
distinto. Descartarlo podría esconder que el archivo trae varios cursos y
violar \ref{req:grado}. Mantenerlo puede, en el peor caso, rechazar un archivo
válido, que es el error barato según la asimetría de costos.

\begin{observacion}[Señal de grados frente a número de páginas]
La señal de que un archivo no es la lista de un grado es el número de grados
distintos, no el número de páginas. En el corpus hay listas de un solo grado
que ocupan seis páginas y archivos con tres grados que caben en dos. Una regla
basada en páginas produce los dos tipos de error a la vez.
\end{observacion}

\subsection{El algoritmo de aceptación}

\Cref{alg:percibir} reúne las piezas. Es una cadena de guardas en un orden
fijo, en el que las comprobaciones baratas van primero: las que no requieren
llamar al modelo preceden a las que sí.

\begin{algorithm}[htbp]
\caption{Percibir un archivo: todo o nada}
\label{alg:percibir}
\begin{algorithmic}[1]
\Require bytes $A$; puertos $\mathrm{Rend}$ (rasterizar) y $\mathrm{Leer}$ (modelo)
\Ensure $(\mathit{ok}, G, P, I, \varepsilon)$
\If{$A$ es vacío} \Return $(0,\varnothing,0,\varnothing,\textsc{ArchivoVacío})$ \EndIf
\State intentar $(I_1,\dots,I_P)\gets\mathrm{Rend}(A)$; si falla
       \Return $(0,\dots,\textsc{PdfIlegible})$
\If{$P > 40$} \Return $(0,\varnothing,P,\varnothing,\textsc{DemasiadoLargo})$
       \Comment{antes de gastar una llamada}
\EndIf
\If{$P = 0$} \Return $(0,\varnothing,0,\varnothing,\textsc{SinPáginas})$ \EndIf
\State $I\gets()$;\quad $\Gamma\gets()$;\quad $\mathit{inc}\gets\varnothing$
\For{$j\gets 1$ \textbf{hasta} $P$}
  \State intentar $(G_j,\kappa_j,\iota_{j,\cdot})\gets\mathrm{Leer}(I_j,j)$
  \State \quad si \textsc{Fallo}: \Return $(0,\varnothing,P,I,\textsc{PercepciónFallida})$
         \Comment{no se devuelve lo leído como lista}
  \If{$\kappa_j = 0$} $\mathit{inc}\gets\mathit{inc}\cup\{j\}$ \EndIf
  \State $\Gamma\gets\Gamma\cdot G_j$
  \State $I\gets I\cdot(\iota_{j,\ell} : \iota_{j,\ell}.x \text{ no es vacío})$
\EndFor
\State $G\gets$ grados distintos de $\Gamma$ bajo $\eta$ \Comment{\cref{def:canon-grado}}
\If{$\card{G} > 1$} \Return $(0,G,P,I,\textsc{VariosGrados})$ \EndIf
\If{$\mathit{inc}\neq\varnothing$} \Return $(0,G,P,I,\textsc{PáginaIncompleta})$ \EndIf
\State \Return $(1,G,P,I,\varnothing)$
\end{algorithmic}
\end{algorithm}

\begin{proposicion}[Corrección del \cref{alg:percibir}]
\label{prop:percibir}
La salida del \cref{alg:percibir} satisface los requisitos
\ref{req:completo}, \ref{req:grado} y \ref{req:nofab}.
\end{proposicion}

\begin{proof}
La única instrucción que devuelve $\mathit{ok}=1$ es la última. Para llegar a
ella, el ciclo debe haber terminado sus $P$ iteraciones sin \textsc{Fallo}
---cualquier fallo retorna dentro del ciclo con $\mathit{ok}=0$---, así que
$I$ es la concatenación de los ítems no vacíos de las $P$ páginas. Además
$\mathit{inc}=\varnothing$, es decir, $\kappa_j=1$ para todo $j$, lo que
prueba \ref{req:completo}. También debe valer $\card{G}\leq1$, por la guarda
anterior, lo que prueba \ref{req:grado}. Los únicos elementos que se agregan a
$I$ provienen de $\mathrm{Leer}$, y $\mathrm{Leer}$ devuelve \textsc{Fallo}, y
no ítems, ante cualquier respuesta que no cumpla el esquema, lo que prueba
\ref{req:nofab}.
\end{proof}

Cuando el archivo se rechaza por varios grados o por una página incompleta,
la salida igual incluye $G$ e $I$. No se usan para armar nada: sirven para
explicarle al acudiente qué se vio («Este archivo tiene la lista de 3 grados:
Cuarto, Quinto, Sexto. Sube solo las páginas del grado de tu hijo»). En un
hogar con varios hijos, la regla se extiende a todo el hogar: si se rechaza el
archivo de un hijo, no se importa el de ninguno.

\paragraph{El tope de páginas.} El límite de 40 páginas no es una señal de
negocio sino una guarda contra archivos abusivos, y se fijó justo por encima
del PDF más largo del corpus (39 páginas, un colegio con todos sus grados en
un solo archivo). Ese archivo sí se lee, y se rechaza por tener trece grados,
diciendo cuáles; esa respuesta le sirve al acudiente. Un archivo más largo que
cualquiera observado se rechaza antes de hacer una sola llamada.

\subsection{Complejidad}

Sea $P$ el número de páginas y $t_j$ el tiempo de la llamada de la página $j$.
La rasterización es lineal en el número total de píxeles,
$\Theta(P\cdot s^2\cdot \mathrm{área})$, y su costo es despreciable frente a
las llamadas. Las llamadas se hacen secuencialmente, así que el tiempo total
es
\[
  T_{\mathrm{per}}(P) = \sum_{j=1}^{P} t_j \;\leq\; P\cdot(2t_{\max}+3),
  \qquad P\leq 40,
\]
con un valor típico de $t_j\approx 5$~s por página. El número de llamadas es
$\Theta(P)$ en el caso sin fallos y como máximo $2P$ con reintentos. La
canonización de grados cuesta $O(\sum_g\card{g})$ y la verificación de
unicidad es lineal en el número de grados mencionados. El costo dominante es
entonces el de $\Theta(P)$ inferencias del modelo, lo que motiva la ruta
rápida: una lista revisada no vuelve a percibirse.

\begin{observacion}[Paralelismo]
Las llamadas son independientes entre páginas, así que podrían hacerse en
paralelo y reducir el tiempo a $\max_j t_j$. La implementación actual es
secuencial, lo que simplifica el orden de los ítems y el control de cuota del
proveedor. Es una mejora directa si la latencia de la primera lectura llega a
ser importante.
\end{observacion}

\subsection{Lo que la percepción no hace}

La percepción \emph{no corrige} lo que lee. Algunos colegios imprimen
«SHSRPIE», «CARPERTAS» o «TAOLLA», y el modelo las transcribe tal cual. Esa
fidelidad es deseable: una corrección hecha dentro del modelo no se puede
auditar ni reproducir. La corrección ortográfica es un algoritmo aparte,
determinista y conservador (\cref{sec:res-ortografia}), que opera contra el
vocabulario del catálogo. Tampoco \emph{clasifica} los renglones ni decide su
tipo: eso corresponde a la normalización (\cref{sec:normalizacion}).

\begin{evolucion}[: percepción]
\begin{itemize}
  \item \textbf{Truncamiento silencioso $\to$ todo o nada.} Un sistema
        anterior limitaba la lectura a las primeras cinco páginas y devolvía el
        resultado como si fuera la lista completa. El efecto era un kit
        incompleto presentado con toda la confianza. La versión vigente
        rechaza el archivo entero cuando supera el tope, y el tope pasó de
        ser un recorte a ser una guarda de 40 páginas, calibrada contra el
        archivo más largo del corpus.
  \item \textbf{Señal por páginas $\to$ señal por grados.} El criterio de
        rechazo pasó a ser el número de grados distintos, con la canonización
        de la \cref{def:canon-grado}.
  \item \textbf{Falsos rechazos por página incompleta.} Dos listas reales
        (Sexto y Décimo de un mismo colegio) se rechazaban porque su segunda
        página terminaba a mitad de la descripción del uniforme, y el modelo
        la marcaba como incompleta. Se agregó la regla S10: un texto que
        continúa en la página siguiente no es una página cortada. La
        protección contra imágenes recortadas o borrosas se mantiene.
  \item \textbf{Segmentación.} Las reglas S3 a S8 se agregaron una por una a
        partir de casos del corpus: contenidos de cartuchera, colores
        enumerados al final, reparto de la cantidad total, alternativas con
        «o», series de durezas y uniformes por nombre. Antes, los uniformes
        salían unas veces prenda por prenda y otras solo con el título. La
        instrucción creció lo mínimo, porque cada línea se paga en cada página
        leída.
  \item \textbf{Reintento único de transporte.} Se agregó tras medir que una
        página colgada 60~s respondía en 6~s al reintentarla. El tiempo límite
        bajó de 60 a 40~s, con lo que el peor caso por página queda en unos
        83~s en lugar de un fallo a los 60.
  \item \textbf{Corrección fuera del modelo.} La corrección ortográfica se
        separó de la percepción para que fuera auditable
        (\cref{sec:res-ortografia}).
\end{itemize}
\end{evolucion}
